%! AP Statistics — Unit 8, Lesson 8.1 — Warm-Up (Answer Key)
\documentclass[10pt]{article}
\usepackage{apstats-article}
\usepackage{apstats-key}

\begin{document}

\pageheader{Unit 8, Lesson 8.1}{Warm-Up --- Answer Key}
\namedateperiod

\vspace{0.15in}

A local coffee shop records the flavor of each drink sold during one week.
The results are shown in the table below.

\vspace{0.1in}
\renewcommand{\arraystretch}{1.4}
\begin{tabularx}{0.55\linewidth}{Xc}
    \textbf{Flavor} & \textbf{Observed Count} \\
    \hline
    Espresso   & 48 \\
    Latte      & 72 \\
    Cappuccino & 36 \\
    Cold Brew  & 44 \\
    \hline
    \textbf{Total} & \textbf{200} \\
\end{tabularx}

\vspace{0.15in}

\begin{enumerate}

  \item What proportion of drinks sold were lattes? Show your work.

        \ansline{$72 \div 200 = 0.36$, or 36\%.}

  \item The shop expected each of the 4 flavors to make up exactly 25\% of total
        sales. How many drinks of each flavor were \emph{expected}?

        \vspace{0.05in}
        \begin{tabularx}{0.8\linewidth}{Xc}
            Expected espresso   & \ans{$0.25 \times 200 = 50$} \\[4pt]
            Expected latte      & \ans{$0.25 \times 200 = 50$} \\[4pt]
            Expected cappuccino & \ans{$0.25 \times 200 = 50$} \\[4pt]
            Expected cold brew  & \ans{$0.25 \times 200 = 50$} \\
        \end{tabularx}

  \item Compare the observed counts to the expected counts. Does the difference
        seem large or small? Explain in 1--2 sentences.

        \ansline{The differences are moderate --- lattes are 22 above expected and cappuccinos are 14 below expected. Whether these differences are due to chance is exactly what chi-square tests will help us determine.}

        \begin{teachernote}
            Accept any reasonable comparison. Students should note that some
            categories are above expected and some are below. The key idea is
            that variation exists and the question is whether it is surprising.
        \end{teachernote}

\end{enumerate}

\end{document}
